% =====================================================================
%  Edward Frenkel, "Langlands Correspondence for Loop Groups"
%  (Cambridge Studies in Advanced Mathematics 103, 2007)
%  中文译注本 — 主控导言区
%  编译：xelatex（两遍；推荐 latexmk -xelatex）
% =====================================================================
\documentclass[11pt,a4paper,UTF8,openany]{ctexbook}

% ---------- 页面与版式 ----------
\usepackage[a4paper,left=2.7cm,right=2.7cm,top=2.7cm,bottom=2.9cm]{geometry}
\usepackage{fancyhdr}
\pagestyle{fancy}
\fancyhf{}
\fancyhead[LE]{\small\itshape Langlands Correspondence for Loop Groups}
\fancyhead[RO]{\small\itshape Langlands Correspondence for Loop Groups}
\fancyhead[LO]{\small\thepage}
\fancyhead[RE]{\small\thepage}
\renewcommand{\headrulewidth}{0.4pt}
\fancypagestyle{plain}{\fancyhf{}\fancyhead[LE,RO]{\small\itshape Langlands Correspondence for Loop Groups}\fancyhead[LO,RE]{\small\thepage}\renewcommand{\headrulewidth}{0.4pt}}

% ---------- 数学 ----------
\usepackage{amsmath}
\usepackage{amssymb}
\usepackage{amsthm}
\usepackage{mathtools}
\usepackage{bm}

% ---------- 图表、颜色、表格、列表 ----------
\usepackage{graphicx}
\usepackage{xcolor}
\usepackage{booktabs}
\usepackage{array}
\usepackage{longtable}
\usepackage{enumitem}
\usepackage{caption}
\usepackage{amscd}
\captionsetup{font=small,labelsep=quad}
\setlength{\emergencystretch}{3em}

% ---------- 超链接 ----------
\usepackage[colorlinks=true,linkcolor=blue!55!black,citecolor=blue!55!black,
            urlcolor=blue!55!black,unicode]{hyperref}

% ---------- 定理类环境 ----------
\theoremstyle{plain}
\newtheorem{theorem}{定理}[section]
\newtheorem{proposition}[theorem]{命题}
\newtheorem{lemma}[theorem]{引理}
\newtheorem{corollary}[theorem]{推论}
\newtheorem{conjecture}[theorem]{猜想}
\theoremstyle{definition}
\newtheorem{definition}{定义}[section]
\newtheorem{example}{例}[section]
\newtheorem{exercise}{习题}[section]
\theoremstyle{remark}
\newtheorem{remark}{注}[section]
\renewcommand{\proofname}{证明}
\renewcommand{\contentsname}{目录}
\renewcommand{\bibname}{参考文献}
\renewcommand{\indexname}{索引}

% ---------- 每节重点（Key Takeaways） ----------
\newenvironment{keypoints}{%
  \par\addvspace{0.9em}%
  \noindent\begin{minipage}{\linewidth}%
  \color{black!80}\rule{\linewidth}{0.6pt}\par\vspace{2pt}%
  \noindent\textbf{【本节重点】}\par\vspace{2pt}%
  \begin{itemize}[leftmargin=1.6em,itemsep=1pt,topsep=2pt,parsep=0pt]%
}{%
  \end{itemize}\vspace{1pt}\noindent\color{black!80}\rule{\linewidth}{0.6pt}%
  \end{minipage}\par\addvspace{0.9em}%
}

% ---------- 正式译注（v2 新增，编译后可见） ----------
% 用法：在正文需要注释的词句后写 \transnote{说明文字 [来源：...]}
% 与源码中的 % [Tutor Note: ...] 并存：后者保留为工作注，前者是给读者的正式译注。
\newcommand{\transnote}[1]{%
  \footnote{\textbf{【译注】}#1}}

% ---------- 常用记号（Frenkel 本书惯例） ----------
\newcommand{\R}{\mathbb{R}}
\newcommand{\C}{\mathbb{C}}
\newcommand{\Z}{\mathbb{Z}}
\newcommand{\Q}{\mathbb{Q}}
\newcommand{\F}{\mathbb{F}}
\newcommand{\PP}{\mathbb{P}}
\newcommand{\GG}{\mathbb{G}}
\newcommand{\dd}{\mathrm{d}}
\newcommand{\ee}{\mathrm{e}}
\newcommand{\ii}{\mathrm{i}}
\newcommand{\Tr}{\operatorname{Tr}}
\newcommand{\tr}{\operatorname{tr}}
\renewcommand{\Re}{\operatorname{Re}}
\renewcommand{\Im}{\operatorname{Im}}
\newcommand{\sgn}{\operatorname{sgn}}
\newcommand{\diag}{\operatorname{diag}}
\newcommand{\id}{\mathrm{id}}
\newcommand{\End}{\operatorname{End}}
\newcommand{\Hom}{\operatorname{Hom}}
\newcommand{\Aut}{\operatorname{Aut}}
\newcommand{\Spec}{\operatorname{Spec}}
\newcommand{\Lie}{\operatorname{Lie}}
\newcommand{\Rep}{\operatorname{Rep}}
\newcommand{\Ind}{\operatorname{Ind}}
\newcommand{\Res}{\operatorname{Res}}
\newcommand{\cO}{\mathcal{O}}
\newcommand{\cA}{\mathcal{A}}
\newcommand{\cB}{\mathcal{B}}
\newcommand{\cC}{\mathcal{C}}
\newcommand{\cD}{\mathcal{D}}
\newcommand{\cE}{\mathcal{E}}
\newcommand{\cF}{\mathcal{F}}
\newcommand{\cG}{\mathcal{G}}
\newcommand{\cH}{\mathcal{H}}
\newcommand{\cI}{\mathcal{I}}
\newcommand{\cJ}{\mathcal{J}}
\newcommand{\cK}{\mathcal{K}}
\newcommand{\cL}{\mathcal{L}}
\newcommand{\cM}{\mathcal{M}}
\newcommand{\cN}{\mathcal{N}}
\newcommand{\cP}{\mathcal{P}}
\newcommand{\cQ}{\mathcal{Q}}
\newcommand{\cR}{\mathcal{R}}
\newcommand{\cS}{\mathcal{S}}
\newcommand{\cT}{\mathcal{T}}
\newcommand{\cU}{\mathcal{U}}
\newcommand{\cV}{\mathcal{V}}
\newcommand{\cW}{\mathcal{W}}
\newcommand{\cX}{\mathcal{X}}
\newcommand{\cY}{\mathcal{Y}}
\newcommand{\cZ}{\mathcal{Z}}
\newcommand{\fg}{\mathfrak{g}}
\newcommand{\fh}{\mathfrak{h}}
\newcommand{\fb}{\mathfrak{b}}
\newcommand{\fn}{\mathfrak{n}}
\newcommand{\fp}{\mathfrak{p}}
\newcommand{\ft}{\mathfrak{t}}
\newcommand{\fu}{\mathfrak{u}}
\newcommand{\fsl}{\mathfrak{sl}}
\newcommand{\fgl}{\mathfrak{gl}}
\newcommand{\fso}{\mathfrak{so}}
\newcommand{\fsp}{\mathfrak{sp}}
\newcommand{\Bun}{\operatorname{Bun}}
\newcommand{\Op}{\operatorname{Op}}
\newcommand{\Conn}{\operatorname{Conn}}
\newcommand{\Loc}{\operatorname{Loc}}
\newcommand{\Gr}{\operatorname{Gr}}
\newcommand{\Fl}{\operatorname{Fl}}
\newcommand{\Ch}{\operatorname{Ch}}
\newcommand{\Crit}{\operatorname{Crit}}
\newcommand{\Sym}{\operatorname{Sym}}
\newcommand{\Tot}{\operatorname{Tot}}
\newcommand{\Coker}{\operatorname{Coker}}
\newcommand{\Image}{\operatorname{Im}}
\newcommand{\gghat}{\hat{\mathfrak{g}}}
\newcommand{\ff}[1]{\mathfrak{#1}}
\newcommand{\norm}[1]{\lVert #1\rVert}
\newcommand{\abs}[1]{\lvert #1\rvert}
